% Einheit 1 von 4 – Die Schar als Familie (bank/scharen-von-geraden-und-ebenen/e1.jsonl, zusammenbau v0.1)
%% TODO Zeitmarke Einheit 1: Marken-Zeile nicht lesbar
%% TODO Zweigzeile Teil 1: was hier gelernt wird (Satz in Schülersprache aus dem Einheitstitel) – Einheit 1
\einheitenkopf[e1]{Einheit 1 von 4 · Die Schar als Familie}
\zweigzeile{Abitur LK}
\setcounter{aufgabe}{8}

%% TODO Titel: Ich-kann-Satz fehlt in der Bank (Platzhalter: Kettenname) – Nr. 9
%% TODO Anweisung: ein Satz über der ersten Teilaufgabe fehlt in der Bank – Nr. 9
\begin{aufgabe}{Was wählt der Parameter?}
\begin{teile}
\teil Wo steht $k$, was verändert es? $g_k: \vec{x} = (2 | 1 | 0) + r \cdot (1 | k | 3)$, $k \in \mathbb{R}$ \\ \kreuz{Stützvektor – die Lage ändert sich, die Richtung bleibt} \\ \kreuz{Richtungsvektor – die Richtung ändert sich} \\ \kreuz{Normalenvektor – die Stellung der Ebene ändert sich} \\ \kreuz{rechte Seite – die Ebenen werden parallel verschoben}
\end{teile}
\end{aufgabe}

%% TODO Titel: Ich-kann-Satz fehlt in der Bank (Platzhalter: Kettenname) – Nr. 10
%% TODO Anweisung: ein Satz über der ersten Teilaufgabe fehlt in der Bank – Nr. 10
\begin{aufgabe}{Die Schar als Familie}
%% TODO \rechenplatz{n}: Schrittzahl des Grundfalls fehlt in der Bank (nur bei fester Schreibform, 2.3 b)
\begin{teile}
\teil Für alle oder für eins? „Zeige, dass $P(1 | 1 | 2)$ in jeder Ebene $E_k: k x_1 - k x_2 + x_3 = 2$ liegt.“ \\ \kreuz{für alle Mitglieder} \\ \kreuz{Parameterwert gesucht}
\teil Zeige, dass $P(2 | 3 | 2)$ in jeder Ebene $E_k: k x_1 + 2x_2 - k x_3 = 6$ liegt ($k \in \mathbb{R}$).
\teil Zeige: $E_5$ schneidet die $x_2$-Achse in $S(0 | 2 | 0)$; $E_k: x_1 + k x_2 + x_3 = 10$.
\teil Gehört $F: 6x_1 + 4x_2 + 8x_3 = 8$ zur Schar $E_k: k x_1 + 2x_2 + (k + 1) x_3 = 4$?
\teil Sind die Geraden $g_k: \vec{x} = (2k | 1 | -k) + r \cdot (1 | 3 | 2)$, $k \in \mathbb{R}$, parallel oder identisch? Begründe. \hfill (Abitur 2024 LK)
\teil Die Punkte $P_k(1 + k | 2 | 3 - 2k)$, $k \in \mathbb{R}$, liegen auf einer Geraden $g$. Gib eine Gleichung von $g$ an. $g: \vec{x} = $ \leerfeld
\teil Begründe mit der Abbildung, dass jede Ebene $E_k: k x_1 + x_2 + 3x_3 = 3$, $k \in \mathbb{R}$, die Gerade $h$ enthält.

\begin{ksys3}\rpunkt{0}{3}{0}{P}\rpunkt{0}{0}{1}{Q}\rgerade[-0.3:1.5]{0,3,0}{0,-3,1}{h}\end{ksys3}
\teil Auf welcher Koordinatenebene stehen alle Ebenen $E_k: k x_1 + 3x_3 = 2k$ senkrecht? \\ \kreuz{$x_1x_2$-Ebene} \\ \kreuz{$x_1x_3$-Ebene} \\ \kreuz{$x_2x_3$-Ebene}
\teil Zeige, dass je zwei verschiedene Ebenen der Schar $E_k: k x_1 + (3 - k) x_2 = 2$, $k \in \mathbb{R}$, nicht parallel sind. \hfill (Abitur 2025 LK)
\end{teile}
\end{aufgabe}

%% TODO Titel: Ich-kann-Satz fehlt in der Bank (Platzhalter: Kettenname) – Nr. 11
%% TODO Anweisung: ein Satz über der ersten Teilaufgabe fehlt in der Bank – Nr. 11
\begin{aufgabe}{Die Schar als Familie – Fehler finden, Begründen}
\begin{teile}
\teil Finde den Fehler. Aufgabe: Zeige, dass $P(2 | 1 | 2)$ in allen Ebenen $E_k: k x_1 + 2x_2 + (1 - k) x_3 = 4$ liegt. \rechnung{k = 0: \quad 0 + 2 + 2 &= 4 \\ k = 1: \quad 2 + 2 + 0 &= 4 && \Rightarrow P \text{ liegt in allen } E_k}
\teil Warum liegt ein Punkt in jeder Ebene einer Schar, wenn beim Einsetzen der Parameter herausfällt?
\end{teile}
\end{aufgabe}
